\chapter{Las primeras fracciones}\label{ch:g4:fractions}

Media manzana, un \omterm{def:g3:division:half}{cuarto} de hora, tres \omterm{def:g3:division:half}{cuartos} de la clase:
las fracciones nombran los trozos de las cosas. Este capítulo presenta la
escritura $\frac{a}{b}$, lee fracciones en dibujos y en la recta numérica
y compara las más sencillas. (Las fracciones crecen en el
\cref{ch:g6:fractions}.)

\section{Nombrar los trozos}

\begin{definition}[Fracción]\label{def:g4:fractions:def}
Corta un todo en partes iguales y toma algunas --- eso es una
\emph{fracción}\index{fracción}:
\[
\frac{3}{4}
\quad
\begin{array}{l}
3 \text{: cuántas partes se toman (el \emph{numerador}\index{numerador});}\\
4 \text{: en cuántas partes iguales se corta el todo (el
\emph{denominador}\index{denominador}).}
\end{array}
\]
$\frac12$ es \emph{un medio}, $\frac13$ \emph{un tercio}, $\frac14$
\emph{un \omterm{def:g3:division:half}{cuarto}} y $\frac{1}{10}$ \emph{un décimo}.
\end{definition}

\begin{omfigure}
\begin{tikzpicture}[scale=0.78]
  % medio circulo
  \draw[thick] (0,0) circle (1.0);
  \fill[omDef!40] (0,0) -- (1,0) arc (0:180:1) -- cycle;
  \draw[thick] (-1,0) -- (1,0);
  \node[below, font=\small] at (0,-1.2) {$\frac12$};
  % cuartos
  \begin{scope}[xshift=3.6cm]
  \draw[thick] (0,0) circle (1.0);
  \fill[omThm!40] (0,0) -- (1,0) arc (0:90:1) -- cycle;
  \fill[omThm!40] (0,0) -- (0,1) arc (90:180:1) -- cycle;
  \fill[omThm!40] (0,0) -- (-1,0) arc (180:270:1) -- cycle;
  \draw[thick] (-1,0) -- (1,0);
  \draw[thick] (0,-1) -- (0,1);
  \node[below, font=\small] at (0,-1.2) {$\frac34$};
  \end{scope}
  % barra de tercios
  \begin{scope}[xshift=6.6cm, yshift=-0.35cm]
  \foreach \i in {0,1,2} \draw[omExo] (\i*1.3,0) rectangle
    (\i*1.3+1.3,0.7);
  \fill[omMeth!40] (0.04,0.04) rectangle (1.26,0.66);
  \draw[thick] (0,0) rectangle (3.9,0.7);
  \node[below, font=\small] at (1.95,-0.15) {$\frac13$};
  \end{scope}
  % barra de decimos
  \begin{scope}[xshift=11.6cm, yshift=-0.35cm]
  \foreach \i in {0,...,9} \draw[omExo] (\i*0.42,0) rectangle
    (\i*0.42+0.42,0.7);
  \foreach \i in {0,...,6} \fill[omProp!40] (\i*0.42+0.03,0.04)
    rectangle (\i*0.42+0.39,0.66);
  \draw[thick] (0,0) rectangle (4.2,0.7);
  \node[below, font=\small] at (2.1,-0.15) {$\frac{7}{10}$};
  \end{scope}
\end{tikzpicture}

{\small Leer fracciones en dibujos. Las partes tienen que ser
\emph{iguales}: ¡tres trozos desiguales no son tercios!}
\end{omfigure}

\begin{example}[Fracciones de una colección]\label{ex:g4:fractions:collection}
$\frac13$ de $12$ canicas: reparte las $12$ canicas en $3$ grupos iguales
de $4$; un grupo es $\frac13$, así que $\frac13$ de $12$ es $4$
--- y $\frac23$ de $12$ son dos grupos, $8$.
\end{example}

\section{Las fracciones en la recta numérica}

\begin{method}[Situar $\frac{a}{b}$]\label{met:g4:fractions:line}
Corta cada unidad de la recta en $b$ pasos iguales; desde $0$, avanza $a$
pasos. Si $a$ es mayor que $b$, el recorrido pasa de $1$: ¡las fracciones pueden
ser mayores que un todo!
\end{method}

\begin{omfigure}
\begin{tikzpicture}[scale=1.6]
  \draw[-Stealth] (-0.2,0) -- (6.9,0);
  \foreach \i in {0,4,8,12,16,20,24}
    \draw ({\i*0.25},-0.09) -- ({\i*0.25},0.09);
  \foreach \i in {0,...,24} \draw ({\i*0.25},-0.05) -- ({\i*0.25},0.05);
  \foreach \x/\l in {0/0, 1/1, 2/2, 3/3, 4/4, 5/5, 6/6}
    \node[below, font=\small] at (\x,-0.1) {$\l$};
  \fill[omDef] (0.5,0) circle (1.6pt)
    node[above, font=\small] {$\frac24$};
  \fill[omThm] (1.25,0) circle (1.6pt)
    node[above, font=\small] {$\frac54$};
  \fill[omMeth] (2.75,0) circle (1.6pt)
    node[above, font=\small] {$\frac{11}{4}$};
\end{tikzpicture}

{\small La recta cortada en \omterm{def:g3:division:half}{cuartos}: $\frac24$ cae en el mismo punto
que $\frac12$; $\frac54$ está un \omterm{def:g3:division:half}{cuarto} pasado el $1$; $\frac{11}{4}$ es
$2$ y $\frac34$.}
\end{omfigure}

\begin{example}[Números enteros escondidos en fracciones]\label{ex:g4:fractions:whole}
$\frac44 = 1$ (cuatro \omterm{def:g3:division:half}{cuartos} hacen el todo); $\frac82 = 4$;
$\frac{12}{3} = 4$. Y $\frac74$ es $1 + \frac34$: un todo y
tres \omterm{def:g3:division:half}{cuartos} más.
\end{example}

\section{Comparar y sumar fracciones sencillas}

\begin{proposition}[Mismo denominador]\label{prop:g4:fractions:compare}
Con el mismo \omterm{def:g4:fractions:def}{denominador}, los trozos tienen el mismo tamaño --- así que basta
con contarlos:
\[
\frac{2}{8} < \frac{5}{8},
\qquad
\frac{3}{8} + \frac{4}{8} = \frac{7}{8}.
\]
Comparar con un todo también es fácil: $\frac{5}{8} < 1 < \frac{9}{8}$
(menos y luego más que $8$ octavos).
\end{proposition}

\begin{proof}[Con un dibujo]
Si sombreas $3$ octavos de una barra y después $4$ octavos más de la misma barra,
quedan sombreados $7$ octavos en total.
\end{proof}

\begin{example}[Denominadores distintos, mismo punto]\label{ex:g4:fractions:equal}
En la recta numérica, $\frac12$, $\frac24$ y $\frac{5}{10}$ caen todas
en el mismo punto: son tres nombres del mismo número. Cortar
cada \omterm{def:g3:division:half}{mitad} en dos da \omterm{def:g3:division:half}{cuartos}; en cinco, décimos. (La regla general
está en el \cref{ch:g6:fractions}.)
\end{example}

\begin{example}[Cuartos de hora]\label{ex:g4:fractions:hour}
Una hora tiene $60$ minutos, así que un \omterm{def:g3:division:half}{cuarto} de hora son
$60 \div 4 = 15$ minutos y tres \omterm{def:g3:division:half}{cuartos} de hora son
$45$ minutos. «Las dos y media» es $\frac12$ hora después de las dos
en punto.
\end{example}

\section{Ejercicios}

\begin{exercise}[$\star$]\label{exo:g4:fractions:1}
Escribe la \omterm{def:g4:fractions:def}{fracción} representada: una pizza cortada en $6$ a la que le quedan $5$ porciones;
una tableta de chocolate de $8$ onzas con $3$ comidas (\omterm{def:g4:fractions:def}{fracción} comida); una
bandera dividida en $3$ franjas verticales iguales con $1$ coloreada.
\end{exercise}

\begin{exercise}[$\star$]\label{exo:g4:fractions:2}
Dibuja una barra cortada en $5$ partes iguales y sombrea $\frac35$. Después sombrea
$\frac{2}{5}$ de otra barra igual. ¿Qué sombreado es mayor?
\end{exercise}

\begin{exercise}[$\star$]\label{exo:g4:fractions:3}
Calcula: $\frac12$ de $18$; \quad $\frac14$ de $20$; \quad
$\frac13$ de $21$; \quad $\frac34$ de $20$ (tres grupos de un
\omterm{def:g3:division:half}{cuarto}).
\end{exercise}

\begin{exercise}[$\star$]\label{exo:g4:fractions:4}
Sitúa en una recta numérica cortada en tercios: $\frac13$; \ $\frac33$; \
$\frac53$; \ $2$. ¿Qué \omterm{def:g4:fractions:def}{fracción} cae exactamente en $2$?
\end{exercise}

\begin{exercise}[$\star$]\label{exo:g4:fractions:5}
Copia y completa con $<$, $>$ o $=$:
\[
\frac38 \;?\; \frac68, \qquad
\frac55 \;?\; 1, \qquad
\frac74 \;?\; 1, \qquad
\frac12 \;?\; \frac24 .
\]
\end{exercise}

\begin{exercise}[$\star$]\label{exo:g4:fractions:6}
Calcula: $\frac26 + \frac36$; \quad $\frac58 - \frac28$; \quad
$\frac14 + \frac24$; \quad $1 - \frac13$ (¿cuántos tercios hacen
$1$?).
\end{exercise}

\begin{exercise}[$\star$]\label{exo:g4:fractions:7}
Escribe como un número entero más una \omterm{def:g4:fractions:def}{fracción} menor que $1$:
$\frac54$; \quad $\frac73$; \quad $\frac{9}{2}$.
\end{exercise}

\begin{exercise}[$\star$]\label{exo:g4:fractions:8}
¿Cuántos minutos son: media hora; un \omterm{def:g3:division:half}{cuarto} de hora; tres
\omterm{def:g3:division:half}{cuartos} de hora; $\frac13$ de hora?
\end{exercise}

\begin{exercise}[$\star$]\label{exo:g4:fractions:9}
En una clase de $24$ alumnos, $\frac14$ llevan gafas. ¿Cuántos
alumnos llevan gafas? ¿Cuántos no?
\end{exercise}

\begin{exercise}[$\star\star$]\label{exo:g4:fractions:10}
Nina se comió $\frac38$ de una pizza y Samuel, $\frac48$ de la misma
pizza. ¿Qué \omterm{def:g4:fractions:def}{fracción} se comieron entre los dos? ¿Qué \omterm{def:g4:fractions:def}{fracción} queda?
¿Quién comió más?
\end{exercise}

\begin{exercise}[$\star\star$]\label{exo:g4:fractions:11}
¿Verdadero o falso? Justifícalo con un dibujo o un ejemplo: «es posible que
$\frac12$ de una pizza pequeña sea menos que $\frac14$ de una pizza muy grande».
¿Qué hay que saber siempre antes de comparar dos fracciones de
algo?
\end{exercise}
